\documentclass[11pt,a4paper]{article}
\usepackage[T1]{fontenc}
\usepackage{lmodern,amsmath,amssymb,amsthm,mathtools,booktabs,microtype}
\usepackage[margin=27mm]{geometry}
\usepackage[colorlinks=true,linkcolor=blue,urlcolor=blue]{hyperref}
\newtheorem{theorem}{Theorem}[section]
\newtheorem{corollary}[theorem]{Corollary}
\newtheorem{remark}[theorem]{Remark}
\newcommand{\Q}{\mathbb Q}\newcommand{\PP}{\mathbb P}
\DeclareMathOperator{\Tr}{Tr}
\setlength{\emergencystretch}{2em}
\title{The surface of ordered matching contact square classes}
\author{}\date{15 November 2026}
\begin{document}\maketitle
\begin{abstract}
Allowing a comparison scale and slope to vary together does not make
square-class matching a rational parametrization problem. For two contact
polynomials of degrees $2g_0+2$ and $2g_1+2$, with simple roots, the smooth
projective model of the ordered correspondence is a surface of general
type when $g_0,g_1\ge2$. Its canonical self-intersection is
$4(g_0-1)(g_1-1)$. This distinguishes the ordered correspondence needed
by the detector arguments from an unordered Kummer quotient. It excludes
a dominant rational parametrization by two free parameters and a
fibration with a rational or genus-one general fiber; it does not exclude
special complete families of arithmetic witnesses.
\end{abstract}

\section{The ordered correspondence}
Let $F_i\in\Q[T]$ be squarefree of even degree $2g_i+2$, with $g_i\ge2$,
for $i=0,1$. Its roots may be rational, but the proof only requires
simplicity over $\overline\Q$. On the open set where $F_0(t_0)F_1(t_1)\ne0$,
the two values have the same class in $\Q^\times/\Q^{\times2}$ exactly when
\[
                         z^2=F_0(t_0)F_1(t_1)
\]
has a rational solution $z$. The two arguments are ordered: interchanging
them is not automatically a permitted identification in the source system.
Homogenizing the $F_i$ gives a double cover of $Y=\PP^1\times\PP^1$ with
branch divisor $D$ consisting of $2g_0+2$ vertical and $2g_1+2$ horizontal
fibers. Write $\pi:X\to Y$ for this normal double cover and
$\widetilde X\to X$ for its minimal resolution.

\begin{theorem}\label{thm:surface}
The surface $\widetilde X$ is smooth, geometrically connected, and of
general type. Its canonical divisor is nef and big, and
\[
                         K_{\widetilde X}^2=4(g_0-1)(g_1-1).
\]
There is no dominant rational map from a rational surface to
$\widetilde X$. No surface birational to $\widetilde X$ admits a
fibration whose smooth general fiber is rational or has genus one.
\end{theorem}
\begin{proof}
A vertical branch valuation of $F_0(t_0)F_1(t_1)$ is odd, so this function
is not a square in $\overline\Q(t_0,t_1)$. The double cover is therefore
geometrically connected. Put
\[
 L=\mathcal O_Y(g_0+1,g_1+1),\qquad D\sim2L.
\]
The cover is smooth above each smooth point of the branch divisor.
There are precisely $(2g_0+2)(2g_1+2)$ crossings. In local coordinates
at a crossing, a unit may be absorbed etale-locally, and the cover has
equation $z^2=uv$, an ordinary double point.

Adjunction for the equation defining a double cover gives
\[
 K_X=\pi^*(K_Y+L)=\pi^*\mathcal O_Y(g_0-1,g_1-1).
\]
This formula also follows directly from the local canonical form
$du\wedge dv/z$ and its changes of coordinates. The resolution of each
ordinary double point is crepant. Indeed in the chart
$u=s$, $z=st$, $v=st^2$ of the resolution the displayed form pulls back to
$2\,ds\wedge dt$, without a zero or pole on the exceptional curve.
The other chart gives the same conclusion. Thus the canonical divisor
of the resolution is the pullback of the divisor just displayed.

The divisor $\mathcal O_Y(g_0-1,g_1-1)$ is ample because both entries are
positive. Its pullback to the finite double cover is ample, and its
pullback under the resolution is nef and big. On $\PP^1\times\PP^1$,
a divisor of bidegree $(a,b)$ has square $2ab$. Multiplication by the
degree two of $\pi$ gives the claimed square of the canonical divisor.
Bigness of the canonical divisor is precisely the asserted general-type
property.

A dominant rational map from a rational surface would be generically
finite. Resolve its indeterminacies. In characteristic zero, pullback
under a generically finite morphism embeds pluricanonical forms of the
target into those of the source (the ramification divisor is effective).
A smooth rational surface has no nonzero pluricanonical forms, whereas
bigness supplies such forms on $\widetilde X$. This is a contradiction.

For the final assertion, resolve a proposed rational fibration on a
smooth surface $Z$ mapping birationally to $\widetilde X$. Let $C$ be
a smooth general fiber and put $K=K_{\widetilde X}$. One has
$C^2=0$ and $K_Z\cdot C=2g(C)-2$. On a succession of point blowups,
$K_Z=f^*K+E$ with $E$ effective and exceptional. Since $C$ is a general
fiber, $E\cdot C\ge0$, and nefness gives $f^*K\cdot C\ge0$.
Genus zero is already impossible. In genus one both intersections
must vanish. As $(f^*K)^2>0$, the Hodge index theorem then implies
$C^2<0$ for the nonzero numerical class of $C$, contrary to $C^2=0$.
Here $C$ is numerically nonzero because it has positive intersection
with an ample divisor on $Z$. This also treats fibrations on other
birational models.
\end{proof}

\section{Consequences for the present source systems}
The replacements $(t_0,t_1)=(\lambda,a\lambda)$ with $\lambda\ne0$,
and $(t_0,t_1)=(\lambda,\lambda+a)$, are birational changes of the
argument plane. Consequently varying $a$ and $\lambda$ together gives
the same surface, not a rational surface obtained by removing the
fixed-scale genus calculation. For equal contact degrees the values are
\[
\begin{array}{c|c|c}
\text{contacts per detector}&g_0=g_1&K_{\widetilde X}^2\\\hline
6&2&4\\
8&3&16\\
10&4&36.
\end{array}
\]
These statements concern matching first contact square classes after
nonzero square factors common to the source have been removed. Any
additional matching of complement classes imposes further conditions.

\begin{remark}[Ordered versus unordered arguments]
When $F_0=F_1=F$, let $C:y^2=F(t)$. The ordered correspondence is
birational to $(C\times C)/\langle(\iota,\iota)\rangle$, where
$\iota$ is the hyperelliptic involution. It is not the additional
quotient by interchanging the two factors. In genus two, that further
unordered quotient is birational to the Kummer quotient of the
Jacobian. Passing to it forgets the ordering and may replace two
rational arguments by a conjugate quadratic pair. Recovering two
\emph{rational} ordered detector arguments is a further condition.
Thus a genus-one fibration on an unordered quotient does not provide
one on the ordered surface in Theorem~\ref{thm:surface}.
\end{remark}

\paragraph{What is not excluded.}
A surface of general type can contain rational curves. In particular,
when $F(2-t)=F(t)$, reflection gives the two rational curves
$t_1=2-t_0$, $z=\pm F(t_0)$. General type does not itself exclude their
use. The separate arithmetic argument in
\nolinkurl{reflection_obstruction_20261114_v32.tex} excludes the actual
reflected eight-contact construction for every sufficiently large
coefficient shift, using the inert prime $11$. Neither theorem excludes
all other matching families, nor proves a variable-count lower bound.
A construction selecting a finite suitable point for each input, or
using a larger nonrational parameter space already present among its
witnesses, is not contradicted by Theorem~\ref{thm:surface}.

\paragraph{Relation to the new quartic.}
The reduction from 204 to 200 variables in
\nolinkurl{quartic200_20261115_v33.tex} avoids this correspondence: its
pole witnesses remain independent. Its saving comes from a recovered
elliptic difference and a shifted integral Lucas equation. No result
here changes the general counts 25 and 13 or the conditional quartic
count 79.

\section*{Verification scope}
The argument in this paper is a mathematical proof, not a finite search
for rational points. The branch count and intersection computation do
not certify the external arithmetic premises of any Diophantine source.
The local double-point calculation is included to make the canonical
divisor calculation auditable. There is no claim of an independent
non-author review.
\end{document}
